\documentclass[11pt]{article}
\usepackage[a4paper,margin=22mm]{geometry}
\usepackage[no-math]{fontspec}
\setmainfont{lmroman10-regular.otf}[BoldFont=lmroman10-bold.otf,ItalicFont=lmroman10-italic.otf,BoldItalicFont=lmroman10-bolditalic.otf]
\usepackage{amsmath,amssymb,amsthm,mathtools,graphicx,xcolor,hyperref,xurl,xspace,ifthen}
\setlength{\emergencystretch}{4em}
\SetMathAlphabet{\mathbf}{normal}{OT1}{cmr}{bx}{n}
\SetMathAlphabet{\mathbf}{bold}{OT1}{cmr}{bx}{n}
\newtheorem{theo}{Theorem}[section]
\newtheorem{lemm}[theo]{Lemma}
\newtheorem{theorem}[theo]{Theorem}
\newtheorem*{remas*}{Remarks}
\newtheorem{remas}[theo]{Remarks}
\newtheorem*{exam*}{Example}
\newtheorem{ques}[theo]{Question}
\providecommand{\og}{«\nobreakspace}
\providecommand{\fg}{\nobreakspace»}
\newtheorem{lem}[theo]{Lemma}
\newtheorem{lemma}[theo]{Lemma}
\newtheorem{prop}[theo]{Proposition}
\newtheorem{coro}[theo]{Corollary}
\newtheorem{corol}[theo]{Corollary}
\newtheorem{conj}[theo]{Conjecture}
\newtheorem{Proposition}[theo]{Proposition}
\newtheorem{theoreme}[theo]{Theorem}
\theoremstyle{definition}
\newtheorem{defi}[theo]{Definition}
\newtheorem{exam}[theo]{Example}
\newtheorem{exem}[theo]{Example}
\newtheorem{remark}[theo]{Remark}
\newtheorem{example}[theo]{Example}
\newtheorem{rema}[theo]{Remark}
\newtheorem{nota}[theo]{Notation}
\newtheorem*{rema*}{Remark}
\newtheorem*{defi*}{Definition}
\newtheorem*{theo*}{Theorem}
\newtheorem*{lemm*}{Lemma}
\newtheorem*{prop*}{Proposition}
\newcommand{\enoncename}{}
\newtheorem{corpusenonce}[theo]{\enoncename}
\newenvironment{enonce}[2][]{\renewcommand{\enoncename}{#2}\begin{corpusenonce}}{\end{corpusenonce}}
\newenvironment{enonce*}[2][]{\par\medskip\noindent\textbf{#2.}\itshape}{\par\medskip}
\newenvironment{proofc}[1][\proofname]{\begin{proof}[#1]}{\end{proof}}
\newcommand{\Mc}{\mathcal{M}}
\newcommand{\pointir}{.\ ---\ }
\providecommand{\mathof}[1]{\mathrm{#1}}

\numberwithin{equation}{section}
\numberwithin{figure}{section}
\numberwithin{table}{section}
\def\endappendix{}
% Add the option ``french'' if the article is in French
% Remove the option ``Unicode'' if your file is latin1-encoded






% Uncomment next line if the editor is a woman
% 





\def\arXiv#1{\href{https://arxiv.org/abs/#1}{arXiv\string:\allowbreak#1}}


%%%%%%%%%%%%MACROS PERSONNELLES%%%%%%%%%%
\usepackage{amsbsy}
\usepackage{amsmath,amsfonts,amssymb,amsthm,mathrsfs,mathtools}
\usepackage{graphicx}
%\usepackage{epsfig}%pour les eps
\usepackage{latexsym}%encore des symboles
\usepackage{xcolor}



\usepackage{xcolor}
%\usepackage[pdfpagemode=UseNone,bookmarksopen=false,colorlinks=true,urlcolor=blue,citecolor=blue,citebordercolor=blue,linkcolor=blue]{hyperref}
%\usepackage{smfhyperref}
\usepackage[capitalize]{cleveref}


%\usepackage{pstricks}
\usepackage{enumerate}
\usepackage{tikz}						%%%%%%%%%%
%\usepackage{boondox}
\usepackage{pifont}
\usepackage{bm}



\usetikzlibrary{arrows, patterns, calc}		%%%%%%%%%%



					%%%%%%%%%%

\DeclareFontFamily{U}{BOONDOX-calo}{\skewchar\font=45 }
\DeclareFontShape{U}{BOONDOX-calo}{m}{n}{
  <-> s*[1.05] BOONDOX-r-calo}{}
\DeclareFontShape{U}{BOONDOX-calo}{b}{n}{
  <-> s*[1.05] BOONDOX-b-calo}{}
\DeclareMathAlphabet{\mathcalbis}{U}{BOONDOX-calo}{m}{n}
\SetMathAlphabet{\mathcalbis}{bold}{U}{BOONDOX-calo}{b}{n}
\DeclareMathAlphabet{\mathbcalboondox}{U}{BOONDOX-calo}{b}{n}

\providecommand{\abs}[1]{\lvert#1\rvert}
\providecommand{\norm}[1]{\lVert#1\rVert}


\newfontfamily\CorpusBbm{Noto Sans Math}
\newcommand{\mathbbm}[1]{\mathord{\text{\CorpusBbm\char"1D7D9}}}
\usepackage{enumerate,booktabs}

\usepackage{enumitem}
	\renewcommand{\theenumi}{{(\roman{enumi})}}
	\renewcommand{\labelenumi}{\theenumi}
	

%%
\renewcommand{\Re}{\operatorname{Re}}
\renewcommand{\Im}{\operatorname{Im}}
%%

\newcommand{\Z}{\mathbb{Z}}
\renewcommand{\P}{\mathbb{P}}
\newcommand{\E}{\mathbb{E}}
\def \NC {\mathbb{NC}}


\def\U{\mathbb{U}}
\def\N{\mathbb{N}}
\newcommand{\R}{\mathbb{R}}
\newcommand{\D}{\mathbb{D}}
\newcommand{\CC}{\mathbb{C}}

\def\S{\mathbb{S}}
\def\e{{\rm e}}
\def\d{{\rm d}}
\def\tx{{\widetilde x}}
\def\eps{\varepsilon}

\def\PPP{\mathcal P}
\def\FFF{\mathcal F}
\def\KKK{\mathcal K}
\def\TTT{\mathcal T}
\def\UUU{\mathcal U}

\def\LLL{\mathcal L}
\def\EEE{\mathcal E}
\def\tmub{\widetilde{\mu}_{\bullet}}
\def\Tree{\mathsf{Tree}}
\newcommand\br[1]{\llbracket #1 \rrbracket}


\def\Vc{\mathcal{V}}
\def\Vcd{\mathcal{V}^{\mathsf{d}}}

%%
\def\i{\mathrm{i}}

%%
\def\FKn{\mathcal{F}^{(n)}_{K_{n}}}
\def\PKn{\mathcal{P}^{(n)}_{K_{n}}}
\def\dFKn{\dot{\mathcal{F}}^{(n)}_{K_{n}}}
\def\dPKn{\dot{\mathcal{P}}^{(n)}_{K_{n}}}
\def\dFKphin{\dot{\mathcal{F}}^{(\phi(n))}_{K_{\phi(n)}}}
\def\dPKphin{\dot{\mathcal{F}}^{(\phi(n))}_{K_{\phi(n)}}}
\def\dFKphin{\dot{\mathcal{P}}^{(\phi(n))}_{K_{\phi(n)}}}
\def\dPKphin{\dot{\mathcal{P}}^{(\phi(n))}_{K_{\phi(n)}}}
\def\dFn{\dot{\mathcal{F}}_{n}}
\def\dPn{\dot{\mathcal{P}}_{n}}
%
\def\ab{a_{\bullet}}
\def\ac{a_{\circ}}
\def\bb{b_{\bullet}}
\def\bc{b_{\circ}}
\def\sc{\sigma_{\circ}}
\def\sb{\sigma_{\bullet}}

%%
\def\abn{a_{\bullet}^{n}}
\def\acn{a_{\circ}^{n}}
\def\bbn{b_{\bullet}^{n}}
\def\bcn{b_{\circ}^{n}}
\def\mub{\mu_{\bullet}}
\def\muc{\mu_{\circ}}

\def\mubn{\mu_{\bullet}^{n}}
\def\mucn{\mu_{\circ}^{n}}

\def\mbn{m_{\bullet}^{n}}
\def\mcn{m_{\circ}^{n}}

\def\scn{\sigma_{\circ}^{n}}
\def\sbn{\sigma_{\bullet}^{n}}

\def\Sbn{S^{\bullet,n}}
\def\Scn{S^{\circ,n}}

\def\phibn{\phi^{\bullet}_{n}}
\def\phicn{\phi^{\circ}_{n}}

\def\DbnN{D^{\bullet}_{n,N}}
\def\Dbn{D^{\bullet}_{n}}

\def\Dcn{D^{\circ}_{n}}

\def\Nb{N^{\bullet}}
\def\Nc{N^{\circ}}

\def\nc{n_{\circ}}
\def\nb{n_{\bullet}}

\def\oB{\overline{B}}
\def\hB{\widehat{B}}
\def\oH{\overline{H}}
\def\hH{\widehat{H}}
\def\oX{\overline{X}}

\def \Xexc{X^{\mathrm{exc}}}
\def \Xbr{X^{\mathrm{br}}}

\def\tTs{\tilde{\Ts}}
\def \Tsb{\Ts^{\bullet}}
\def \Tsc{\Ts^{\circ}}
\def \Tsbn{\Ts^{\bullet,n}}
\def \Tscn{\Ts^{\circ,n}}
\newcommand\BGW{\textrm{BGW}}

%%
\def\tb{\mathcalbis{t}}
\def\Ts{\mathscr{T}}
\renewcommand{\epsilon}{\varepsilon}

\newcommand\Es[1]{\mathbb{E}\left[#1\right]}
\newcommand\Esb[1]{\mathbb{E}\big[#1\big]}
\newcommand\EsB[1]{\mathbb{E}\Big[#1\Big]}
\renewcommand\Pr[1]{\mathbb{P}\left(#1\right)}
\newcommand\tPr[1]{\tilde{\mathbb{P}}\left(#1\right)}
\newcommand\Prb[1]{\mathbb{P}\big(#1\big)}
\newcommand\PrBr[1]{\mathbb{P}\Big(#1\Big)}
\newcommand\PrInd[2]{\mathbb{P}_#1\left(#2\right)}
\newcommand \fl[1] {\left\lfloor #1 \right\rfloor}
\newcommand \ce[1] {\left\lceil #1 \right\rceil}
%\newcommand{\comment}[1]{\todo[color=green!40,inline]{#1}}
%\newcommand{\igor}[1]{\todo[color=blue!20,inline]{I, #1}}

%%


\def\llbracket{[\hspace{-.10em} [ }
\def\rrbracket{ ] \hspace{-.10em}]}
\def\f{\mathcal{F}}
\def\ve{\varepsilon}
\def\la{{\longrightarrow}}
\def\build#1_#2^#3{\mathrel{
\mathop{\kern 0pt#1}\limits_{#2}^{#3}}}
\def\wt{\widetilde}
\def\One{\mathbbm{1}}

\DeclareMathOperator{\Find}{Find}
\DeclareMathOperator{\id}{id}


%%%%%%%%%%FIN MACROS PERSONNELLES%%%%%%%%

\graphicspath{{./figures/}}

%\urladdr{}









\begin{document}
\section*{\raggedright 2746: Uniform inverse Gaussian local limits for Cayley-weighted offspring sums when the white population is of square-root order}
\noindent\textbf{Author:} Valentin Féray; Igor Kortchemski\par\noindent\textbf{Source:} The geometry of random minimal factorizations of a long cycle via biconditioned bitype random trees. Annales Henri Lebesgue 1 (2018), publisher version, DOI 10.5802/ahl.5. \url{https://ahl.centre-mersenne.org/item/AHL_2018__1__149_0/}
\par\noindent\textbf{License:} CC BY 4.0, \url{https://creativecommons.org/licenses/by/4.0/}. The original publisher PDF has the explicit license notice. See the saved source/license records.
\par\noindent\textbf{Editorial material note (not author text):} The selected problem is Lemma 3.10 for the exact Cayley-weighted triangular offspring laws, Kn/sqrt(n) tending to c>0, fixed u>0, the original integer-summand condition and the authored |j|<=n\textasciicircum{}(3/8) uniformity. The original notation table, complete parameter/variance lemma, both local-limit statements and complete local-limit proof subsection retain every local new core, including the moving square-root singularity, intermediate-frequency bound and aperiodic-frequency estimate. Standard Lambert expansion, singular differentiation, inverse-Gaussian characteristic function and Fourier inversion remain named prerequisites. The black local-limit proof is retained because the author explicitly refers to its inversion steps. Tree coding, functional bridges, laminations, figures and animation are unselected and unnecessary. All proof text and formulas below are editable author TeX. Mathematical wording, language and proof order are retained. Only the article class, title matter and bibliography presentation are adapted for independent compilation; no proof has been generated or repaired.
\par\noindent\textbf{Editorial difficulty:} H1: A triangular offspring law approaches its square-root generating-function singularity while the number of summands grows only at square-root scale. The author supplies a uniform local limit through separate moving-singularity, intermediate-frequency and aperiodic-frequency estimates, beyond the fixed-law central/local limit machinery of ordinary qualifiers.
\par\noindent\textbf{Editorial retrieval and attribution:} Retrieved 2026-10-01. DOI 10.5802/ahl.5. Source attribution notice: ../licenses/2746-source-license.txt; complete legal text: ../licenses/CC-BY-4.0.txt. Original source excerpts and the existing source-material check are retained in ../sources/2746/.
\medskip\hrule\medskip
\noindent\textbf{Author statement, complete proof and local supporting text follow in source order.}
\setcounter{section}{2}
\section{Limit theorems for bi-conditioned bi-type BGW trees}
\label{sec:limitheorems}


\begin{table}[htbp]\caption{Table of the main notation and symbols appearing in \cref{sec:limitheorems}.}
\centering
\begin{tabular}{c c p{0.75 \linewidth} }
\toprule
$X, \Xbr,\Xexc$ && The L\'evy process characterized by (3.1),  its associated  bridge process going from $0$ to $0$ in a unit time and its associated excursion process.\\
$\mubn$, $\mucn$&& The offspring distributions defined by \eqref{eq:mubn_mucn}.\\
$\mathscr{T}_{n}$ & & An alternating two-type BGW tree, with offspring distributions $\mubn$ and $\mucn$ conditioned on having $n-K_{n}$ black vertices and $K_{n}+1$ white vertices.
 \\
 $\mathscr{T}^{n}$ & & An alternating two-type BGW tree, with offspring distributions $\mubn$, $\mucn$.\\
 $\mbn,(\sbn)^{2}$ & & Resp.~the mean and variance of $\mubn$.\\
 $ \mcn,(\scn)^{2}$ & & Resp.~the mean and variance of $\mucn$.\\
 $\Scn_{k}$ , $\Sbn_{k}$ && Resp.~the sum of $k$ i.i.d.\ random variables distributed according to $\mucn,\mubn$.\\
\bottomrule
\end{tabular}
\label{tab:seclim}
\end{table}
\setcounter{subsection}{2}\setcounter{theo}{4}\setcounter{equation}{9}
\subsection{A functional invariance principle for alternating BGW trees}
\label{ss:inv}

It is now time to consider specific offspring distributions.
These distributions will be chosen of the form
\begin{equation}
  \mubn(i)= \abn \cdot  (\bbn)^{i} \cdot \frac{(i+1)^{i-1}}{i!}  \quad (i \geq 0), \qquad \mucn(i)= \acn \cdot  (\bcn)^{i} \cdot \frac{(i+1)^{i-1}}{i!}  \quad (i \geq 0).
  \label{eq:mubn_mucn}
\end{equation}
Indeed, according to Proposition~2.5,
alternating BGW trees with these distributions are connected to minimal factorizations
(for the moment, let us forget that, in Proposition~2.5, the root has a different offspring distribution).
We start by specifying the choices of the parameters
$\abn$, $\bbn$, $\acn$ and $\bcn$.
(Note that the chosen parameters and hence the offspring distributions depend on $n$.)

\begin{lemma}\label{lem:exist} For every $n \ge 1$,  fix $K_{n} \in \{1,\dots,n-1\}$.
\begin{enumerate}
  \item[(i)] We may choose positive parameters  $\abn,\bbn,\acn,\bcn$ in a unique way  such that \eqref{eq:mubn_mucn}
defines probability distributions with respective means
\[m^{n}_{\bullet}  \, \coloneqq  \, \sum_{i=0}^{\infty} i \cdot \mubn(i)= \frac{K_{n}+1}{n-K_{n}}, \qquad m^{n}_{\circ} \, \coloneqq \, \sum_{i=0}^{\infty} i \cdot \mucn(i)= \frac{n-K_{n}}{K_{n}+1}.\]
\item[(ii)] Assume that $K_{n} \rightarrow \infty$. If $\tfrac{K_{n}}{n} \rightarrow 0$ as $n \rightarrow \infty$, then
\[(\sbn)^{2} \sim \frac{K_{n}}{n}, \quad  \bbn \sim \frac{K_{n}}{n}  \qquad \text{ and } \qquad  (\scn)^{2} \sim   \big( \frac{n}{K_{n}} \big)^{3},\]
where $(\sbn)^{2}$ and $(\scn)^{2}$ denote respectively the variance of $\mubn$ and $\mucn$.
In particular, if $c>0$ is fixed and $\tfrac{K_{n}}{\sqrt{n}} \rightarrow c$ as $n \rightarrow \infty$, then   $ (\sbn)^{2} \sim \tfrac{c}{\sqrt{n}}$ and $(\scn)^{2} \sim  \frac{n^{3/2}}{c^{3}}$.
\end{enumerate}
\end{lemma}

\begin{proof}Let $F$ be the power series defined by
\begin{equation}
\label{eq:F}F(z)= \sum_{k \geq 0} \frac{(k+1)^{k-1}}{k!} z^{k}.
\end{equation}
It is elementary to check that $G(z)= \frac{z F'(z)}{F(z)}$ defines a continuous increasing function on $[0,1/e)$ and that $\lim_{z \rightarrow 0+} G(z)=0$, $\lim_{z \rightarrow 1/e-} G(z)=\infty$. One then chooses $\bbn, \bcn$ such that respectively $G(\bbn)=\frac{K_{n}+1}{n-K_{n}}$, $G(\bcn)= \frac{n-K_{n}}{K_{n}+1}$, and then $\abn=1/F(\bbn)$, $\acn=1/F(\bcn)$.
This proves (i).

For (ii), since $G(z)=z+o(z)$ as $z \rightarrow 0$ and since $G(\bbn)= \frac{K_{n}+1}{n-K_{n}}$, we get that $\bbn \sim \frac{K_{n}}{n}$ as $n \rightarrow \infty$.
The estimate for the variance $(\sbn)^{2}$ is then obtained by using the expression
\hbox{$(\sbn)^{2}= \frac{(\bbn)^{2} F''(\bbn)}{F(\bbn)}+\mbn-(\mbn)^{2}$:}
the dominant term when $n  \rightarrow \infty$ is $\mbn$
and we have \hbox{$(\sbn)^{2}\sim \frac{K_{n}}{n}$}.


The estimate concerning $(\scn)^{2}$ is more subtle, as it involves the behavior of $F$ near its radius of convergence $1/e$.
We can analytically extend $F$
on a complex split neighborhood on $1/e$ (i.e.\ on a complex neighborhood of $1/e$, without the real half-line $[1/e,+\infty)$).
Indeed we have $F(z)=-W(-z)/z$, where $W$ is the Lambert function (see \cite[Eq~3.1]{CGHJK96}).
Using \cite[Eq~4.22]{CGHJK96}, as $z \rightarrow 0$, we have
\begin{equation}
\label{eq:devF}F \left( \frac{1}{e}-z \right)=e- \sqrt{2} e^{3/2} \sqrt{z}+ o(\sqrt{z}),
\end{equation}
where $z$ is a complex number avoiding the negative real line
(throughout the paper, we use the principal determination of
$\sqrt{z}$ when $z$ is in $\CC \backslash \R_{<0}$).
By singular differentiation \cite[Theorem VI.8 p. 419]{FS09},
we get expansions for $F' \big( \tfrac{1}{e}-z \big)$ and $F'' \big( \tfrac{1}{e}-z \big)$
by differentiating the right-hand side of \eqref{eq:devF}.
Hence $G(1/e-z) \sim \frac{1}{\sqrt{2e}} \cdot \frac{1}{\sqrt{z}}$ as $z \rightarrow 0$.
By definition $\bcn$ is the solution of $G(\bcn)=\tfrac{n-K_n}{K_n+1}$, so that
\begin{equation}
  1/e-\bcn  \quad \mathop{\sim}_{n \rightarrow \infty} \quad  \frac{1}{2e} \left( \frac{K_n}{n-K_n}  \right)^{2}  \quad \mathop{\sim}_{n \rightarrow \infty} \quad  \frac{1}{2e} \left( \frac{K_n}{n}  \right)^{2}.
  \label{eq:est_bcirc}
\end{equation}
We again use the expression \hbox{$(\scn)^{2}= \frac{(\bcn)^{2} F''(\bcn)}{F(\bcn)}+m_{\circ}-m_{\circ}^{2}$}
to estimate the variance.
An easy computation gives $(\scn)^{2}\sim ( {n}/{K_{n}} )^{3}$.
(This time, the dominant term is $\left({(\bcn)^{2} F''(\bcn)}\right)\big/{F(\bcn)}$.)
\end{proof}
\setcounter{theo}{8}\setcounter{equation}{14}
\begin{lemma} \label{lem:ll1}
  Fix $0<u \leq 1$. For $n,N \geq 1$,  set  $\DbnN= \sqrt{ (\sbn)^{2}  N}$. Assume that $ \tfrac{K_{n}}{n} \rightarrow 0$ and that $K_{n} \rightarrow \infty$ as $n \rightarrow \infty$.
  We have
  \[  \sup_{un \leq N  \leq n} \sup_{k \in \Z  } \left|  \DbnN  \cdot \Pr{\Sbn_N=k}
  -   p \left( \frac{k- N\tfrac{K_n+1}{n-K_n}}
  {  \DbnN  }\right)\right|  \quad \mathop{\longrightarrow}_{n \rightarrow \infty} \quad 0,\]
  where $p(x)= \frac{1}{\sqrt{2 \pi}} e^{- \frac{x^{2}}{2}}$ is the standard Gaussian density.
  \label{lem:ll1Uniforme}
\end{lemma}



\begin{lemma}\label{lem:ll2}
 Fix $u>0$. Assume that $ \tfrac{K_{n}}{\sqrt{n}} \rightarrow c>0$ as $n \rightarrow \infty$.  We set $\Dcn=\sqrt{ (\scn)^{2} K_{n} } \sim n /c$.
 Then it holds that
\[ \sup_{|j| \leq n^{3/8}} \sup_{k \in \Z} \left|  \Dcn \cdot \Pr{\Scn_{ uK_{n}  + j }=k} -  q_{u} \left( \frac{k}{  \Dcn }\right)\right|  \quad \mathop{\longrightarrow}_{n \rightarrow \infty} \quad 0,\]
where
\begin{equation}
\label{eq:defqu}q_{u}(x)=\left( \frac{u^{2} c^{3}}{2 \pi x^{3}} \right)^{1/2} \exp \left(  - \frac{c (x-uc)^{2} }{2x} \right) \mathbbm{1}_{x>0}.
\end{equation}\end{lemma}


We have chosen the exponent $3/8 \in (1/4,1/2)$ in view of specific future use,  but as the proof will show we can replace $3/8$ by  any positive value less than $1/2$.

\setcounter{subsection}{4}\setcounter{theo}{15}\setcounter{equation}{30}
\subsection{Local limit theorems}
\label{ss:llt}

We now establish two local limit estimates concerning the random walks $\Sbn$ and $\Scn$
(Lemmas~\ref{lem:ll1} and \ref{lem:ll2}).



\paragraph*{Local limit theorem for $\Sbn$.}


Let $\phibn(t)=\Esb{e^{it\Sbn_{1}}}= \tfrac{F( \bbn e^{it})}{F(\bbn)}$  be the characteristic function of $\Sbn_{1}$,
i.e.\ of the distribution $\mubn$.
(Recall that $F(z)=\sum_{k \ge 0} \tfrac{(k+1)^{k-1}}{k!} z^k$ was introduced in the proof of Lemma~\ref{lem:exist}.)

\begin{lemma}\label{lem:phib}
Assume that $\tfrac{K_{n}}{n} \rightarrow 0$ and $K_{n} \rightarrow \infty$ as $n \rightarrow \infty$. For every $n$ sufficiently large, for every $ 0 \leq |t| \leq \pi$  we have
\begin{equation}
\label{eq:phib}  \ln | \phibn (t)| \leq -  \frac{K_{n}}{ 8n} t^{2}.
\end{equation}
\end{lemma}

\begin{proof}
We take two different approaches depending on the distance of $|t|$ to $0$.

First assume that $|t| \leq 1$. We have (see e.g.~\cite[Lemma 3.3.7]{Dur10})
\[\phibn(t)=1+\textrm{i}\E[\Sbn_{1}]t- \E[(\Sbn_{1})^{2}] \frac{t^{2}}{2}+r^{n}_{1}(t),\]
with $|r^{n}_{1}(t)| \leq |t|^{3} \E[(\Sbn_{1})^{3}]/6$ for every $n \geq 1$ and $t \in \R$. By using the expansion of $F$ around $0$, we see that $ \E[\Sbn_{1}] \sim \E[(\Sbn_{1})^{2}] \sim \E[(\Sbn_{1})^{3}] \sim \frac{K_{n}}{n} \rightarrow 0$ as $n \rightarrow \infty$, so that
\begin{align*}
  \hspace{ -2mm} |\phibn(t)|^{2} &= \left( 1+\textrm{i}\E[\Sbn_{1}]t- \E[(\Sbn_{1})^{2}] \frac{t^{2}}{2}+r^{n}_{1}(t) \right) \left( 1-\textrm{i}\E[\Sbn_{1}]t- \E[(\Sbn_{1})^{2}] \frac{t^{2}}{2}+\overline{r^{n}_{1}(t)} \right)\\
&= 1 - (\sbn)^{2} t^{2}+r^{n}_{2}(t),
\end{align*}
where $|r^{n}_{2}(t)| \leq  |t|^{3} K_{n}/(2n)$ for every $n$ sufficiently large, uniformly for $t \in [-1,1]$.
Therefore
\[ 2 \ln |\phibn(t)| \leq - (\sbn)^{2} t^{2}+ |t|^{3} K_{n}/(2n)= - \frac{K_{n}}{n} t^{2} \left(  \frac{(\sbn)^{2} n}{K_{n}}-|t|/2 \right)
\]
for every $n$ sufficiently large, uniformly for $t \in [-1,1]$.
Since $(\sbn)^{2} \sim K_{n}/n$ as $n \rightarrow \infty$ by Lemma~\ref{lem:exist} (ii), it follows that for every $n$ sufficiently large and $0 \leq |t| \leq 1$ we have $ \ln | \phibn (t)| \leq -  \frac{K_{n}}{ 8n} t^{2}$.

Now assume that $1 \leq |t| \leq \pi$. Since $F(z)=1+z+o(z)$ as $z \rightarrow 0$ and, by definition, $\phibn(t)=\tfrac{F( \bbn e^{it})}{F(\bbn)}$,
we have
\[|\phibn(t)|^{2}=(1+\bbn(e^{it}-1)+o(\bbn))(1+\bbn(e^{-it}-1)+o(\bbn))=1+2\bbn(\cos(t)-1)+o(\bbn),\]
where the $o(\bbn)$ is uniform in $1\leq |t| \leq \pi$.
Therefore, for every $1 \leq |t| \leq \pi$, for every $n$ sufficiently large,
\[ 2 \ln |\phibn(t)| \leq - \bbn \frac{t^{2}}{3}+o(\bbn) \leq - \frac{K_{n}}{4n} t^{2},\]
where for the last inequality we have used the fact that $\bbn \sim \frac{K_{n}}{n}$ as $n \rightarrow \infty$ by Lemma~\ref{lem:exist} (ii) and that $t$ is not too close to $0$.
\end{proof}





\begin{proof}[Proof of Lemma~\ref{lem:ll1}]
We first check that the convergence
\begin{equation}
\label{eq:cv1b}\E\bigg[e^{i t \frac{\Sbn_{N}-N \frac{K_{n}+1}{n-K_{n}} }{ \DbnN }} \bigg]  \quad \mathop{\longrightarrow}_{n \rightarrow \infty} \quad e^{- \frac{t^{2}}{2}}
\end{equation}
holds uniformly for $un \leq N \leq n$ and $t$ in compact subsets of $\R$.  As in the proof of Lemma~\ref{lem:phib}, since $\tfrac{\E[(\Sbn_{1})^{3}]}{(\DbnN)^{3}}=\mathcal O(\tfrac{1}{n \sqrt{K_{n}}})= o(\frac{1}{n})$, we have
\[\phibn\left( \frac{t}{ \DbnN } \right)=1+\textrm{i}\E[\Sbn_{1}] \frac{t}{ \DbnN } - \E[(\Sbn_{1})^{2}] \frac{t^{2}}{2(\DbnN)^{2}}+ o \left( \frac{1}{n} \right),\]
where the $o$ is uniform when $un \leq N \leq n$ and $t$ belongs to a compact subset of $\R$, so that
\[\ln \phibn \left( \frac{t}{ \DbnN } \right)=it  \frac{K_{n}+1}{(n-K_{n})\DbnN} - \frac{t^{2}}{2} \cdot \frac{(\sbn)^{2}}{(\DbnN )^{2}}+ o \left( \frac{1}{n} \right).\]
Thus
\[\Es{e^{i t \frac{\Sbn_{N}-N \frac{K_{n}+1}{n-K_{n}} }{ \DbnN }}}= \exp \left(- \frac{t^{2}}{2} \cdot  \frac{ N (\sbn)^{2} }{(\DbnN)^{2}} + o \left( \frac{N}{n} \right)   \right).\]
Since $(\DbnN)^{2}= N (\sbn)^{2}$, this  implies \eqref{eq:cv1b}.


We then follow the steps of the analytic proof of the standard local limit theorem for a sequence of independent identically distributed random variables. The main difficulty is that the distribution of $\Sbn_{1}$ depends on $n$. To simplify notation, set $f(t)= e^{- \frac{t^{2}}{2}}$ for $t \in \R$. By Fourier inversion we have, for every $x \in \mathbb{R}$,

\begin{equation}
 \label{eq:fourierb}  p(x)= \frac{1}{2\pi} \int_{-\infty}^{\infty} e^{-itx} f(t) {\d}t.
 \end{equation}
Also, for $k \in \mathbb{Z}$,
\[ \Pr{\Sbn_{N}=k} = \frac{1}{2\pi} \int_{-\pi}^{\pi} e^{-itk} \phibn(t)^{ N} {\mathrm{d}}t.\]
In the following, we only consider values of $x \in \R$ such that $N \frac{K_{n}+1}{n-K_{n}}+x  \DbnN$ is an integer.
For such $x$,
\begin{multline*}
  \DbnN \cdot \Pr{\Sbn_{ N}= N \frac{K_{n}+1}{n-K_{n}}+x  \DbnN}\\
  =  \frac{1}{2\pi} \int^{\pi \DbnN}_{-\pi \DbnN} e^{- i t x} \left(  \phibn \left(  \frac{t}{\DbnN} \right)   e^{-it \frac{K_{n}+1}{(n-K_{n}) \DbnN}} \right)^{N} {\mathrm{d}}t.
\end{multline*}
Therefore for every fixed $A>0$, for $n$ sufficiently large,
\begin{multline*}
  \left| \DbnN \cdot \Pr{\Sbn_{N}= N \frac{K_{n}+1}{n-K_{n}}+x  \DbnN} -p( x)\right| \\
  \leq \frac{1}{2 \pi} \big(|I^{(1)}_{A}(x,N,n)| + |I^{(2)}_{A}(x,N,n)|+|I^{(3)}(x)|\big),
\end{multline*}
where
\begin{align*}
  I^{(1)}_{A}(x,N,n)&=\int_{-A}^{A} e^{-i t x} \left(   \left(  \phibn \left(  \frac{t}{\DbnN} \right)   e^{-it \frac{K_{n}+1}{(n-K_{n}) \DbnN}} \right)^{N}- f(t) \right) {\mathrm{d}}t,\\
 I^{(2)}_{A}(x,N,n)&= \int_{A<|t|< \pi \DbnN} e^{-itx - it \frac{(K_{n}+1) N}{(n-K_{n}) \DbnN}} \phibn \left(  \frac{t}{\DbnN} \right) ^{ N}{\mathrm{d}}t, \\
 I^{(3)}_{A}(x)&= \int_{|t|>A} e^{-itx} f(t) {\mathrm{d}}t.
 \end{align*}
We shall check that for every fixed $\epsilon>0$, there exists $A>0$ such that for every $n$ sufficiently large, for every $ un \leq N \leq  n$, for every $x \in \R$ (with the above integrality condition),  $|I^{(1)}_{A}(x,N,n)| \leq \epsilon$, $|I^{(2)}_{A}(x,N,n)| \leq \epsilon$ and $|I^{(3)}_{A}(x)| \leq\epsilon$.



 We choose $A >0$ such that
\begin{equation}
\label{eq:Ab} 2 \int_{A}^{\infty } e^{-  \frac{1}{9} t^{2}} {\d}t < \epsilon.
\end{equation}


\emph{Bounding $|I^{(1)}_{A}(x,N,n)|$.} Since the convergence \eqref{eq:cv1b} holds  uniformly on compact subsets of $\R$, for   $n$ sufficiently large, for every $ un \leq N \leq n$ and $x \in \R$, we have $|I^{(1)}_{A}(x,N,n)| \leq \epsilon$.

\emph{Bounding $|I^{(3)}_{A}(x)|$.} By the choice of $A$ in \eqref{eq:Ab}
since $|f(t)|=e^{-  \frac{1}{2} t^{2}} \le e^{-  \frac{1}{9} t^{2}}$,
for every $x \in \R$ we have $|I^{(3)}_{A}(x)| \leq \epsilon$.


 \emph{Bounding $|I^{(2)}_{{A} }(x,N,n)|$.}  By Lemma~\ref{lem:phib},  for every $n$ sufficiently large, for every $ un \leq N \leq n$, $A<|t| \leq \pi  \DbnN$ we have
\[|I^{(2)}(x,N,n)| \leq 2 \int_{A}^{\pi \DbnN} e^{-  \frac{ K_{n} N }{ 8 n (\DbnN)^{2}} \cdot t^{2} } {\d}t \leq 2 \int_{A}^{\infty } e^{-  \frac{1}{9} t^{2}} {\d}t,\]
   which is less than $\epsilon$ by \eqref{eq:Ab}. This completes the proof.
\end{proof}





\paragraph*{Local limit theorem for $\Scn$.}

In the case of $\Scn$, the analysis is more subtle.
As above, we start with an estimate on the characteristic function
of the step distribution of the random walk
\[\phicn(t)\coloneqq\Es{e^{it\Scn_{1}}}=\frac{F( \bcn e^{it})}{F(\bcn)}.\]
To simplify notation, set  $y_{n}=1/e-\bcn$.

\begin{lemma}\label{lem:technique}
Assume that $\tfrac{K_{n}}{\sqrt{n}} \rightarrow c>0$ as $n \rightarrow \infty$.  There exist constants $A_{0},\epsilon,\kappa>0$, which may only depend on $c$, such that the following holds.
For every $n$ sufficiently large and for every $|t| \in (\frac{A_{0}}{\Dcn},\epsilon)$, we have $|\phicn(t)| \leq e^{- \kappa |t|^{1/2}}$.
\end{lemma}



\begin{proof}We start with some preliminary observations which fix the values of the different constants.
\begin{enumerate}
\item[--]  Since $\Dcn \sim n /c$ and $y_{n} \sim \frac{c^{2}}{2en}$ (as was seen in the proof of Lemma~\ref{lem:exist}),
  we can find a constant $\gamma>0$ depending only on $c$ such that
  the conditions $t \ge  {A_0}/{\Dcn}$ and $A_{0} \geq 1$ imply
  $t \ge {\gamma A_0}/{n}$ and $t \ge \gamma y_{n}$,
  for $n$ sufficiently large.
\item[--] We then fix $A_{0} \geq 1$ such that $\frac{2c}{\sqrt{\gamma A_{0}}} \leq  \frac{1}{4\sqrt{e}}$.
\item[--] From the expansion $ \ln F \big( \tfrac{1}{e}-z \big)= 1-\sqrt{2e} \sqrt{z}+o(\sqrt{z})$, there exists $\eta>0$ such that
\begin{equation}
\label{eq:z} \Re(z) \ge 0, \, |z| \leq \eta  \qquad \implies  \qquad  \Re \ln F \left( \frac{1}{e}-z \right) \leq  1-  \Re \sqrt{z}.
\end{equation}
Then set $\epsilon= \sqrt{\eta^{2}/(e^{-2}+\gamma^{-2})}$.
\end{enumerate}


We now turn to the main part of the proof. For $t \in \R$, since
$ \ln |\phicn(t)| =\Re \ln \phicn(t)$, we have
\begin{equation}
\label{eq:lnphi} \ln |\phicn(t)| =\Re \ln F( \bcn e^{it})- \ln F( \bcn).
\end{equation}
We first study the term $\ln F( \bcn)$.
By \eqref{eq:devF} and \eqref{eq:est_bcirc}, we have $\ln F( \bcn)=1- \tfrac{c}{\sqrt{n}}+ o \big(  \frac{1}{\sqrt{n}} \big)$.  As a consequence, for every $n$ sufficiently large, for every $  {A_{0}}/{\Dcn} \leq |t| \leq 1$, we have
\begin{equation}
\label{eq:ln}\ln F( \bcn) \geq 1- \frac{2c}{\sqrt{n}} \geq 1- \frac{2c}{\sqrt{\gamma A_{0}}} \sqrt{|t|}.
\end{equation}

Consider now the term $\Re \ln F( \bcn e^{it})$.
To this end, we set $z=z_{n}(t)=\tfrac{1}{e}-\bcn e^{it}$.
Note that we always have $\Re(z_n(t)) \ge 0$.
Observing that $e^{-it} z_{n}(t)= \tfrac{1}{e}(e^{-it}-1)+y_n$, we have
\begin{equation}
  \label{eq:module} |z_{n}(t)|^{2}=  \tfrac{1}{e^2} \big|e^{-it}-1\big|^2
  +\tfrac{2}{e} \big\langle y_n, e^{-it}-1 \big\rangle +|y_n|^2
  = \frac{2(1-\cos(t))}{e^{2}} -\frac{2 y_{n}(1-\cos(t))}{e}+y_{n}^{2}.
\end{equation}
For $n$ sufficiently large and $|t| \leq \epsilon$, it follows that
\begin{equation}
\label{ineq:module} |z_{n}(t)|^{2} \leq  \left( \frac{1}{e^{2}} + \frac{1}{\gamma^{2}} \right)  t^{2} \leq \eta^{2}.
\end{equation}
Using \eqref{eq:z}, we have
\[\Re \ln F( \bcn e^{it}) \leq 1-   \Re \sqrt{z_{n}(t)}.\]
Since $\Re z_{n}(t) \ge 0$, we have
\[ \Re \sqrt{z_{n}(t)}=\sqrt{\frac{\Re(z_n(t))+|z_n(t)|}{2}} \ge \sqrt{\frac{|z_n(t)|}{2}}. \]
But, by \eqref{eq:module}, for $n$ sufficiently large and  for every $|t|$ in $(A_{0}/\Dcn,1)$,
\[|z_{n}(t)|^{2} = \frac{2(1-\cos(t))}{e^{2}} (1-y_{n} e)+y_{n}^{2} \geq \frac{1-\cos(t)}{e^{2}} \geq  \frac{t^{2}}{4e^{2}}.\]
As a consequence $\Re \sqrt{z_{n}(t)} \geq \frac{\sqrt{|t|}}{2\sqrt{e}}$, so that
\begin{equation}
\label{eq:reln}\Re \ln F( \bcn e^{it}) \leq 1- \frac{\sqrt{|t|}}{2\sqrt{e}}.
\end{equation}
To conclude, by combining \eqref{eq:lnphi} with \eqref{eq:ln} and \eqref{eq:reln},
we get that for every $n$ sufficiently large, for every $|t|$ in $(A_{0}/\Dcn, \epsilon)$,
\[\ln|\phicn(t)| \leq \left(  \frac{2c}{\sqrt{\gamma A_{0}}}- \frac{1}{2\sqrt{e}} \right) \cdot \sqrt{|t|} \leq -\frac{1}{4\sqrt{e}} \sqrt{|t|}.\]
This completes the proof, with $\kappa= \frac{1}{4\sqrt{e}}$.
\end{proof}


We are now ready to tackle the proof of Lemma~\ref{lem:ll2}.
\begin{proof}[Proof of Lemma~\ref{lem:ll2}]
The first step is to show the convergence
\begin{equation}
\label{eq:cv1}\Es{e^{i t \frac{\Scn_{ u K_{n} +j}}{ \Dcn }}}  \quad \mathop{\longrightarrow}_{n \rightarrow \infty} \quad \exp \left( u c^{2} \left( 1 -\sqrt{1-  \frac{2it}{c}}  \right)  \right),
\end{equation}
uniformly for $t$ in compact subsets of $\R$ and $|j| \leq n^{3/8}$.

The characteristic function $\phicn(t)=\Es{e^{it\Scn_{1}}}$  of $\Scn_{1}$
is given as $\phicn(t)= \frac{F( \bcn e^{it})}{F(\bcn)}$,
where $\bcn$ has been chosen such that $ \frac{\bcn F'(\bcn)}{F(\bcn)}= \frac{n-K_{n}}{K_{n}+1}$.
Recall that $\bcn=1/e-y_{n}$ with
\hbox{$y_{n} \sim \frac{1}{2e}  \big( \frac{K_{n}}{n} \big)^{2} \sim \frac{c^{2}}{2e} \cdot \frac{1}{n}$}
(from \eqref{eq:est_bcirc}).
Since $ \Dcn \sim  n /c$, we have, uniformly on compact subsets of $\R$ as $n \rightarrow \infty$,
\[\bcn e^{ \frac{it}{\Dcn}}=(1/e-y_{n})e^{ \frac{it}{\Dcn}}=1/e-  \left( \frac{c^{2}}{2e} - \frac{it c}{e } \right) \cdot \frac{1}{{n}}+o \left(  \frac{1}{n} \right).\]
By \eqref{eq:devF}, we have $ \ln F \left( \frac{1}{e}-z \right)= 1-\sqrt{2e} \sqrt{z}+o(\sqrt{z})$ as $z \rightarrow 0$.
As a consequence as $n \rightarrow \infty$, \[\ln \phicn \left( \frac{t}{ \Dcn } \right)=-\sqrt{2e} \sqrt{ \frac{c^{2}}{2e} - \frac{it c}{e}} \cdot \frac{1}{\sqrt{n}} + \sqrt{2e} \sqrt{ \frac{c^{2}}{2e}} \cdot \frac{1}{\sqrt{n}}+o \left(  \frac{1}{\sqrt{n}} \right),\]
 uniformly on compact subsets of $\R$. Therefore
\[( uK_{n}  +j) \ln \phicn \left( \frac{t}{\Dcn} \right)  \quad \mathop{\longrightarrow}_{n \rightarrow \infty} \quad u c^{2}-uc \sqrt{c^{2}- 2itc},\]
 uniformly when $t$ belongs to compact subsets of $\R$ and $|j| \leq n^{3/8}$, which implies \eqref{eq:cv1}.

 As in the proof of Lemma~\ref{lem:ll1}, the second step consists in estimating $\Pr{\Scn_{ uK_{n}  +j}=k}$ by Fourier inversion.
 We first let, for $t \in \mathbb{R}$,
\[f(t)=\exp \left( u c^{2} \left( 1 -\sqrt{1-  \frac{2it}{c}}  \right)  \right)\]
be the expression appearing in \eqref{eq:cv1}.
By \cite[Sec.1.2.5]{Kyp06}),  $f$ is the characteristic function of a random variable with explicit density
\[q_u(x)=\left( \frac{u^{2} c^{3}}{2 \pi x^{3}} \right)^{1/2} \exp \left(  - \frac{c (x-uc)^{2} }{2x} \right) \mathbbm{1}_{x>0},\]
so that by Fourier inversion we have, for every $x \in \mathbb{R}$,
\begin{equation}
 \label{eq:fourier}
 q_u(x)=  \frac{1}{2\pi}\int_{-\infty}^{\infty}  {\mathrm{d}t} e^{-itx} \exp \left( u c^{2} \left( 1 -\sqrt{1-  \frac{2it}{c}}  \right)  \right).
 \end{equation}
On the other hand,
for $k \in \mathbb{Z}$,
\[ \Pr{\Scn_{ uK_{n}  +j}=k} = \frac{1}{2\pi} \int_{-\pi}^{\pi} e^{-itk} \phicn(t)^{ uK_{n}  +j} {\mathrm{d}}t.\]
Therefore, assuming that $x \in \R$ is chosen so that $x \Dcn$ is an integer, we get that
\[ \Dcn \cdot \Pr{\Scn_{ uK_{n}  +j}=x  \Dcn}=  \frac{1}{2\pi} \int^{\pi \Dcn}_{-\pi \Dcn} e^{- i t x} \phicn \left(  \frac{t}{\Dcn} \right) ^{ uK_{n}  +j} {\mathrm{d}}t.\]
We deduce that for every fixed $A>0$ and $0< \epsilon \leq 1$, for $n$ sufficiently large, it holds
\begin{multline*}
  \left| \Dcn \cdot \Pr{\Scn_{ uK_{n}  +j}=x  \Dcn} - q_{u}(x)\right| \\
  \leq \frac{1}{2 \pi} (|I^{(1)}_{A}(x,j,n)| + |I^{(2)}_{\epsilon,A}(x,j,n)|+|I^{(3)}_{\epsilon}(x,j,n)|+|I^{(4)}_{A}(x)|),
\end{multline*}
where
\[I^{(1)}_{A}(x,j,n)=\int_{-A}^{A} e^{-i t x} \left(  \phicn \left(  \frac{t}{\Dcn} \right) ^{ uK_{n}  +j}- f(t) \right) {\mathrm{d}}t\]
\[ I^{(2)}_{\epsilon,A}(x,j,n)= \int_{A<|t|< \epsilon \Dcn} e^{-itx} \phicn \left(  \frac{t}{\Dcn} \right) ^{ uK_{n}  +j}{\mathrm{d}}t,\]
\[ I^{(3)}_{\epsilon}(x,j,n)=\int_{\epsilon \Dcn<|t|<\pi \Dcn} e^{-itx} \phicn  \left(  \frac{t}{\Dcn} \right) ^{ uK_{n}  +j} {\mathrm{d}}t, \qquad I^{(4)}_{A}(x)= \int_{|t|>A} e^{-itx} f(t) {\mathrm{d}}t .\]
We shall check that for every fixed $\epsilon'>0$, there exists $A>0$ and $0 < \epsilon \leq 1$ such that for every $n$ sufficiently large, for every $|j| \leq n^{3/8}$, for every $x \in \R$,
 $|I^{(1)}_{A}(x,j,n)| \leq \epsilon', |I^{(2)}_{\epsilon,A}(x,j,n)| \leq \epsilon', |I^{(3)}_{\epsilon}(x,j,n)| \leq \epsilon', |I^{(4)}_{A}(x)| \leq\epsilon'$.


Fix $\epsilon'>0$. We first explain how to choose $A$ and $\epsilon$.
As a preliminary observation, note that there exists a constant $C_{1}>0$
(which only depend on $u$ and $c$) such that $|f(t)| \leq C_{1} e^{-C_{1}^{-1} \sqrt{t}}$ for $t \geq 0$
(this is straightforward, using the explicit expression of $f$).
Let $A_{0}$ and $\kappa$ be the constants given by Lemma~\ref{lem:technique}.
We may choose $A \geq A_{0}$ such that
\begin{equation}
\label{eq:A}2\int_{A}^{\infty} |f(t)|  \mathrm{d}t <\epsilon' \qquad \textrm{and} \qquad  2 \int_{A}^{\infty} e^{- \frac{\kappa u c^{3/2}}{2} t^{1/2}} \mathrm{d}t < \epsilon'.
\end{equation}
We then let $\epsilon$ be also given by Lemma~\ref{lem:technique}.






\emph{Bounding $|I^{(1)}_{A}(x,j,n)|$.} Since the convergence \eqref{eq:cv1} holds  uniformly on compact subsets of $\R$ and $|j| \leq n^{3/8}$, for   $n$ sufficiently large, for every $|j| \leq n^{3/8}$ and $x \in \R$, we have $|I^{(1)}_{A}(x,j,n)| \leq \epsilon'$.

\emph{Bounding $|I^{(4)}_{A}(x)|$.} By the choice of $A$ in \eqref{eq:A}, for every $n$ sufficiently large, for every $x \in \R$ we have $|I^{(4)}_{A}(x,n)| \leq \epsilon'$.

\emph{Bounding $|I^{(3)}_{\epsilon}(x,j,n)|$.}  Set $\phi(t)=\frac{F(e^{it}/e)}{F(1/e)}$, which is the characteristic function of  the probability distribution $ \frac{(i+1)^{(i-1)}}{e^{i+1} i!}$ on $\Z_{+}$. Since the latter is non lattice, we have $|\phi(t)|<1$ for $t \in (0,2\pi)$ (see e.g.~\cite[Theorem 3.5.1]{Dur10}).
Therefore, there exists $C_2>0$ such that $|\phi(t)|<e^{-2C_2}$ for $ \epsilon \leq |t| \leq \pi$.
Since $F$ is a power series with nonnegative coefficients and converges in $1/e$,
it converges uniformly on $\{z: |z| \le 1/e\}$ and is continuous on this compact set.
It is therefore uniformly continuous and
we have $\phicn(t) \rightarrow \phi(t)$, uniformly for $t \in \R$.
In particular, for $n$ sufficiently large and $t$ with $ \epsilon \leq |t| \leq \pi$,
we have $|\phicn(t)|<e^{-C_2}$.
 Therefore, for every $n$ sufficiently large, for every $|j| \leq n^{3/8}$ and  $x \in \R$,
 \[ \left|I^{(3)}_{\epsilon}(x,j,n) \right|  \leq 2 \int_{\epsilon \Dcn}^{\pi \Dcn} e^{-C_2 (uK_{n}  +n^{3/8})} \mathrm{d}t
 \leq 2 \pi \Dcn \cdot (e^{- C_2 u c \sqrt{n}/2} ) \le \epsilon'.\]

 \emph{Bounding $|I^{(2)}_{\epsilon, {A} }(x,j,n)|$.}  By Lemma~\ref{lem:technique}, for every $A<|t| \leq \epsilon \Dcn$ we have \[  \left| \phicn \left(  \frac{t}{\Dcn} \right) \right|  \leq e^{-\kappa \left(  \frac{|t|}{\Dcn} \right) ^{1/2}}.\]
Hence, for $n$ sufficiently large, for every $|j| \leq n^{3/8}$ and $x \in \R$,
 \[ \left| I^{(2)}_{\epsilon,A}(x,j,n)  \right| \leq 2 \int_{A}^{\epsilon \Dcn}\left( e^{-\kappa  \left(  \frac{t}{\Dcn} \right) ^{1/2}} \right) ^{ uK_{n}  -n^{3/8}} \mathrm{d}t .\]
 Since $K_{n} \sim c \sqrt{n}$ and $ \Dcn\sim  n /c$, it follows that for $n$ sufficiently large, for every $|j| \leq n^{3/8}$ and $x \in \R$,
  \[ \left| I^{(2)}_{\epsilon,A}(x,j,n) \right| \leq 2 \int_{A}^{\epsilon \Dcn}  e^{- \frac{\kappa u  c^{3/2}}{2} t^{1/2}} \mathrm{d}t  \leq 2 \int_{A}^{\infty} e^{- \frac{\kappa u c^{3/2}}{2} t^{1/2}} \mathrm{d}t,\]
  which is less than $\epsilon'$ by \eqref{eq:A}.   This completes the proof.
  \end{proof}
\begin{thebibliography}{MMMMMM}
\bibitem[ACEH16]{ACEH16} Alexandrov, Alexander; Chapuy, Guillaume; Eynard, Bertrand; Harnad, John Weighted Hurwitz numbers and topological recursion: an overview (2016) (arXiv:1610.09408) \href{https://ahl.centre-mersenne.org/item/AHL_2018__1__149_0/#ACEH16}{Publisher reference}.
\bibitem[AHRV07]{AHRV07} Angel, Omer; Holroyd, Alexander E.; Romik, Dan; Virág, Bálint Random sorting networks, Adv. Math., Volume 215 (2007) no. 2, pp. 839-868 \href{https://ahl.centre-mersenne.org/item/AHL_2018__1__149_0/#AHRV07}{Publisher reference}.
\bibitem[AL11]{AL11} Armendáriz, Inés; Loulakis, Michail Conditional distribution of heavy tailed random variables on large deviations of their sum, Stochastic Process. Appl., Volume 121 (2011) no. 5, pp. 1138-1147 \href{https://ahl.centre-mersenne.org/item/AHL_2018__1__149_0/#AL11}{Publisher reference}.
\bibitem[Ald94a]{Ald94a} Aldous, David Recursive self-similarity for random trees, random triangulations and Brownian excursion, Ann. Probab., Volume 22 (1994) no. 2, pp. 527-545 \href{https://ahl.centre-mersenne.org/item/AHL_2018__1__149_0/#Ald94a}{Publisher reference}.
\bibitem[Ald94b]{Ald94b} Aldous, David Triangulating the circle, at random, Amer. Math. Monthly, Volume 101 (1994) no. 3, pp. 223-233 \href{https://ahl.centre-mersenne.org/item/AHL_2018__1__149_0/#Ald94b}{Publisher reference}.
\bibitem[AP98]{AP98} Aldous, David; Pitman, Jim The standard additive coalescent, Ann. Probab., Volume 26 (1998) no. 4, pp. 1703-1726 \href{https://ahl.centre-mersenne.org/item/AHL_2018__1__149_0/#AP98}{Publisher reference}.
\bibitem[BD06]{BD06} Berestycki, Nathanaël; Durrett, Rick A phase transition in the random transposition random walk, Probab. Theory Related Fields, Volume 136 (2006) no. 2, pp. 203-233 \href{https://ahl.centre-mersenne.org/item/AHL_2018__1__149_0/#BD06}{Publisher reference}.
\bibitem[Ber96]{Ber96} Bertoin, Jean Lévy processes, Cambridge Tracts in Mathematics, 121, Cambridge University Press, 1996 \href{https://ahl.centre-mersenne.org/item/AHL_2018__1__149_0/#Ber96}{Publisher reference}.
\bibitem[Ber11]{Ber11} Berestycki, Nathanaël Emergence of giant cycles and slowdown transition in random transpositions and k-cycles, Electron. J. Probab., Volume 16 (2011) no. 5, pp. 152-173 \href{https://ahl.centre-mersenne.org/item/AHL_2018__1__149_0/#Ber11}{Publisher reference}.
\bibitem[Ber18]{Ber16} Berzunza, Gabriel On scaling limits of multitype Galton-Watson trees with possibly infinite variance, ALEA Lat. Am. J. Probab. Math. Stat., Volume 15 (2018) no. 1, pp. 21-48 \href{https://ahl.centre-mersenne.org/item/AHL_2018__1__149_0/#Ber16}{Publisher reference}.
\bibitem[Bet18]{Bet17} Bettinelli, Jérémie Convergence of uniform noncrossing partitions toward the Brownian triangulation, Sém. Lotharigien Combin., Volume 80B (2018), \#38, 12 pages (FPSAC Proceedings) \href{https://ahl.centre-mersenne.org/item/AHL_2018__1__149_0/#Bet17}{Publisher reference}.
\bibitem[BFG04]{BDFG04} Bouttier, Jérémie; Francesco, Philippe Di; Guitter, Emmanuel Planar maps as labeled mobiles, Electron. J. Combin., Volume 11 (2004) no. 1 Research Paper 69, 27 pp. (electronic) \href{https://ahl.centre-mersenne.org/item/AHL_2018__1__149_0/#BDFG04}{Publisher reference}.
\bibitem[Bia97]{Biane1997} Biane, Philippe Some properties of crossings and partitions, Discrete Math., Volume 175 (1997) no. 1-3, pp. 41-53 \href{https://ahl.centre-mersenne.org/item/AHL_2018__1__149_0/#Biane1997}{Publisher reference}.
\bibitem[Bia02]{Biane2002} Biane, Philippe Parking functions of types a and b, Electron. J. Combin., Volume 9 (2002) no. 1 (Note \#7, 5 pp) \href{https://ahl.centre-mersenne.org/item/AHL_2018__1__149_0/#Biane2002}{Publisher reference}.
\bibitem[Bia05]{Biane2004} Biane, Philippe Nombre de factorisations d’un grand cycle, Sém. Lothar. Combin., Volume 51 (2004/05) (Art. B51a, 4 pp) \href{https://ahl.centre-mersenne.org/item/AHL_2018__1__149_0/#Biane2004}{Publisher reference}.
\bibitem[Bil68]{Bil68} Billingsley, Patrick Convergence of probability measures, first edition, Wiley, New-York, 1968 \href{https://ahl.centre-mersenne.org/item/AHL_2018__1__149_0/#Bil68}{Publisher reference}.
\bibitem[BK00]{BK00} Bennies, Jürgen; Kersting, Götz A random walk approach to Galton-Watson trees, J. Theoret. Probab., Volume 13 (2000), pp. 777-803 \href{https://ahl.centre-mersenne.org/item/AHL_2018__1__149_0/#BK00}{Publisher reference}.
\bibitem[Bra14]{UB14} Bravo, Gerónimo Uribe Bridges of Lévy processes conditioned to stay positive, Bernoulli, Volume 20 (2014) no. 1, pp. 190-206 \href{https://ahl.centre-mersenne.org/item/AHL_2018__1__149_0/#UB14}{Publisher reference}.
\bibitem[CB11]{CUB11} Chaumont, Loïc; Bravo, Gerónimo Uribe Markovian bridges: weak continuity and pathwise constructions, Ann. Probab., Volume 39 (2011) no. 2, pp. 609-647 \href{https://ahl.centre-mersenne.org/item/AHL_2018__1__149_0/#CUB11}{Publisher reference}.
\bibitem[CG11]{CLGrecursive} Curien, Nicolas; Gall, Jean-François Le Random recursive triangulations of the disk via fragmentation theory, Ann. Probab., Volume 39 (2011) no. 6, pp. 2224-2270 \href{https://ahl.centre-mersenne.org/item/AHL_2018__1__149_0/#CLGrecursive}{Publisher reference}.
\bibitem[CGH + 96]{CGHJK96} Corless, R. M.; Gonnet, G. H.; Hare, D. E. G.; Jeffrey, D. J.; Knuth, D. E. On the Lambert W function, Adv. Comput. Math., Volume 5 (1996) no. 4, pp. 329-359 \href{https://ahl.centre-mersenne.org/item/AHL_2018__1__149_0/#CGHJK96}{Publisher reference}.
\bibitem[CK14]{CK14} Curien, Nicolas; Kortchemski, Igor Random non-crossing plane configurations: a conditioned Galton-Watson tree approach, Random Structures Algorithms, Volume 45 (2014) no. 2, pp. 236-260 \href{https://ahl.centre-mersenne.org/item/AHL_2018__1__149_0/#CK14}{Publisher reference}.
\bibitem[CK15]{CK15} Curien, Nicolas; Kortchemski, Igor Percolation on random triangulations and stable looptrees, Probab. Theory Related Fields, Volume 163 (2015) no. 1-2, pp. 303-337 \href{https://ahl.centre-mersenne.org/item/AHL_2018__1__149_0/#CK15}{Publisher reference}.
\bibitem[CL16]{CL16} Chaumont, Loïc; Liu, Rongli Coding multitype forests: application to the law of the total population of branching forests, Trans. Amer. Math. Soc., Volume 368 (2016) no. 4, pp. 2723-2747 \href{https://ahl.centre-mersenne.org/item/AHL_2018__1__149_0/#CL16}{Publisher reference}.
\bibitem[CW13]{CWmht} Curien, Nicolas; Werner, Wendelin The Markovian hyperbolic triangulation, J. Eur. Math. Soc., Volume 15 (2013) no. 4, pp. 1309-1341 \href{https://ahl.centre-mersenne.org/item/AHL_2018__1__149_0/#CWmht}{Publisher reference}.
\bibitem[Dau18]{dauvergne2018archimedean} Dauvergne, Duncan The archimedean limit of random sorting networks (2018) (arXiv:1802.08934) \href{https://ahl.centre-mersenne.org/item/AHL_2018__1__149_0/#dauvergne2018archimedean}{Publisher reference}.
\bibitem[DMWZZ04]{DMZZ04} Diaconis, Persi; Mayer-Wolf, Eddy; Zeitouni, Ofer; Zerner, Martin P. W. The Poisson-Dirichlet law is the unique invariant distribution for uniform split-merge transformations, Ann. Probab., Volume 32 (2004) no. 1B, pp. 915-938 \href{https://ahl.centre-mersenne.org/item/AHL_2018__1__149_0/#DMZZ04}{Publisher reference}.
\bibitem[DS81]{DS81} Diaconis, Persi; Shahshahani, Mehrdad Generating a random permutation with random transpositions, Z. Wahrsch. Verw. Gebiete, Volume 57 (1981) no. 2, pp. 159-179 \href{https://ahl.centre-mersenne.org/item/AHL_2018__1__149_0/#DS81}{Publisher reference}.
\bibitem[Dur10]{Dur10} Durrett, Rick Probability: theory and examples, Cambridge Series in Statistical and Probabilistic Mathematics, 31, Cambridge University Press, Cambridge, 2010 \href{https://ahl.centre-mersenne.org/item/AHL_2018__1__149_0/#Dur10}{Publisher reference}.
\bibitem[Dén59]{Den59} Dénes, József The representation of a permutation as the product of a minimal number of transpositions, and its connection with the theory of graphs, Magyar Tud. Akad. Mat. Kutató Int. Közl., Volume 4 (1959), pp. 63-71 \href{https://ahl.centre-mersenne.org/item/AHL_2018__1__149_0/#Den59}{Publisher reference}.
\bibitem[EG87]{EG87} Edelman, Paul; Greene, Curtis Balanced tableaux, Adv. Math., Volume 63 (1987) no. 1, pp. 42-99 \href{https://ahl.centre-mersenne.org/item/AHL_2018__1__149_0/#EG87}{Publisher reference}.
\bibitem[ELSV01]{ELSV} Ekedahl, Torsten; Lando, Sergei; Shapiro, Michael; Vainshtein, Alek Hurwitz numbers and intersections on moduli spaces of curves, Invent. Math., Volume 146 (2001) no. 2, pp. 297-327 \href{https://ahl.centre-mersenne.org/item/AHL_2018__1__149_0/#ELSV}{Publisher reference}.
\bibitem[FK18]{FK18} Féray, Valentin; Kortchemski, Igor Trajectories in random minimal transposition factorizations (2018) (arXiv:1810.07586) \href{https://ahl.centre-mersenne.org/item/AHL_2018__1__149_0/#FK18}{Publisher reference}.
\bibitem[FS09]{FS09} Flajolet, Philippe; Sedgewick, Robert Analytic combinatorics, Cambridge University Press, Cambridge, 2009 \href{https://ahl.centre-mersenne.org/item/AHL_2018__1__149_0/#FS09}{Publisher reference}.
\bibitem[Fér12]{Fer12} Féray, Valentin Partial Jucys-Murphy elements and star factorizations, Eur. J. Combin., Volume 33 (2012), pp. 189-198 \href{https://ahl.centre-mersenne.org/item/AHL_2018__1__149_0/#Fer12}{Publisher reference}.
\bibitem[Gal05]{LG05} Gall, Jean-François Le Random trees and applications, Probability Surveys (2005), pp. 245-311 \href{https://ahl.centre-mersenne.org/item/AHL_2018__1__149_0/#LG05}{Publisher reference}.
\bibitem[GJ99a]{GJ99B} Goulden, I. P.; Jackson, Daniel M. The number of ramified coverings of the sphere by the double torus, and a general form for higher genera, J. Combin. Th. Ser. A, Volume 88 (1999) no. 2, pp. 259-275 \href{https://ahl.centre-mersenne.org/item/AHL_2018__1__149_0/#GJ99B}{Publisher reference}.
\bibitem[GJ99b]{GJ99A} Goulden, I. P.; Jackson, Daniel M. A proof of a conjecture for the number of ramified coverings of the sphere by the torus, J. Combin. Th. Ser. A, Volume 88 (1999) no. 2, pp. 246-258 \href{https://ahl.centre-mersenne.org/item/AHL_2018__1__149_0/#GJ99A}{Publisher reference}.
\bibitem[GJ09]{GouldenJackson2009} Goulden, I. P.; Jackson, Daniel M. Transitive powers of Young-Jucys-Murphy elements are central, Journal of Algebra, Volume 321 (2009) no. 7, pp. 1826-1835 \href{https://ahl.centre-mersenne.org/item/AHL_2018__1__149_0/#GouldenJackson2009}{Publisher reference}.
\bibitem[GM11]{LGM09} Gall, Jean-François Le; Miermont, Grégory Scaling limits of random planar maps with large faces, Ann. Probab., Volume 39 (2011) no. 1, pp. 1-69 \href{https://ahl.centre-mersenne.org/item/AHL_2018__1__149_0/#LGM09}{Publisher reference}.
\bibitem[GP93]{GP93} Goulden, I. P.; Pepper, S. Labelled trees and factorizations of a cycle into transpositions, Discrete Math., Volume 113 (1993) no. 1-3, pp. 263-268 \href{https://ahl.centre-mersenne.org/item/AHL_2018__1__149_0/#GP93}{Publisher reference}.
\bibitem[GP08]{LGP08} Gall, Jean-François Le; Paulin, Frédéric Scaling limits of bipartite planar maps are homeomorphic to the 2-sphere, Geom. Funct. Anal., Volume 18 (2008) no. 3, pp. 893-918 \href{https://ahl.centre-mersenne.org/item/AHL_2018__1__149_0/#LGP08}{Publisher reference}.
\bibitem[GY02]{GY02} Goulden, Ian; Yong, Alexander Tree-like properties of cycle factorizations, J. Combin. Theory Ser. A, Volume 98 (2002) no. 1, pp. 106-117 \href{https://ahl.centre-mersenne.org/item/AHL_2018__1__149_0/#GY02}{Publisher reference}.
\bibitem[Hoe63]{Hoeffding} Hoeffding, Wassily Probability inequalities for sums of bounded random variables, J. Amer. Stat. Assoc., Volume 58 (1963) no. 301, pp. 13-30 \href{https://ahl.centre-mersenne.org/item/AHL_2018__1__149_0/#Hoeffding}{Publisher reference}.
\bibitem[Hur91]{Hur91} Hurwitz, Adolf Ueber Riemann’sche Flächen mit gegebenen Verzweigungspunkten, Math. Ann., Volume 39 (1891) no. 1, pp. 1-60 \href{https://ahl.centre-mersenne.org/item/AHL_2018__1__149_0/#Hur91}{Publisher reference}.
\bibitem[IR09]{IrvingRattan2009} Irving, John; Rattan, Amarpreet Minimal factorizations of permutations into star transpositions, Discrete Mathematics, Volume 309 (2009) no. 6, pp. 1435-1442 \href{https://ahl.centre-mersenne.org/item/AHL_2018__1__149_0/#IrvingRattan2009}{Publisher reference}.
\bibitem[Jac88]{Jac88} Jackson, Daniel Some combinatorial problems associated with products of conjugacy classes of the symmetric group, J. Combin. Theory Ser. A, Volume 49 (1988), pp. 363-369 \href{https://ahl.centre-mersenne.org/item/AHL_2018__1__149_0/#Jac88}{Publisher reference}.
\bibitem[Jan12]{Jan12} Janson, Svante Simply generated trees, conditioned Galton-Watson trees, random allocations and condensation, Probab. Surv., Volume 9 (2012), pp. 103-252 \href{https://ahl.centre-mersenne.org/item/AHL_2018__1__149_0/#Jan12}{Publisher reference}.
\bibitem[JS03]{JS03} Jacod, Jean; Shiryaev, Albert N. Limit theorems for stochastic processes, Grundlehren der Mathematischen Wissenschaften, 288, Springer-Verlag, Berlin, 2003 \href{https://ahl.centre-mersenne.org/item/AHL_2018__1__149_0/#JS03}{Publisher reference}.
\bibitem[Kal81]{Kal81} Kallenberg, Olav Splitting at backward times in regenerative sets, Ann. Probab., Volume 9 (1981) no. 5, pp. 781-799 \href{https://ahl.centre-mersenne.org/item/AHL_2018__1__149_0/#Kal81}{Publisher reference}.
\bibitem[Kal02]{Kal02} Kallenberg, Olav Foundations of modern probability, Probability and its Applications, Springer-Verlag, New York, 2002 \href{https://ahl.centre-mersenne.org/item/AHL_2018__1__149_0/#Kal02}{Publisher reference}.
\bibitem[KM16]{KM16} Kortchemski, Igor; Marzouk, Cyril Triangulating stable laminations, Electron. J. Probab., Volume 21 (2016) (Paper No. 11, 31 p.) \href{https://ahl.centre-mersenne.org/item/AHL_2018__1__149_0/#KM16}{Publisher reference}.
\bibitem[KM17]{KM17} Kortchemski, Igor; Marzouk, Cyril Simply Generated Non-Crossing Partitions, Combin. Probab. Comput., Volume 26 (2017) no. 4, pp. 560-592 \href{https://ahl.centre-mersenne.org/item/AHL_2018__1__149_0/#KM17}{Publisher reference}.
\bibitem[Kor14]{Kor14} Kortchemski, Igor Random stable laminations of the disk, Ann. Probab., Volume 42 (2014) no. 2, pp. 725-759 \href{https://ahl.centre-mersenne.org/item/AHL_2018__1__149_0/#Kor14}{Publisher reference}.
\bibitem[Kor15]{Kor15} Kortchemski, Igor Limit theorems for conditioned non-generic Galton-Watson trees, Ann. Inst. Henri Poincaré Probab. Stat., Volume 51 (2015) no. 2, pp. 489-511 \href{https://ahl.centre-mersenne.org/item/AHL_2018__1__149_0/#Kor15}{Publisher reference}.
\bibitem[Kyp06]{Kyp06} Kyprianou, Andreas E. Introductory lectures on fluctuations of Lévy processes with applications, Universitext, Springer-Verlag, Berlin, 2006 \href{https://ahl.centre-mersenne.org/item/AHL_2018__1__149_0/#Kyp06}{Publisher reference}.
\bibitem[Lin92]{Lin92} Lindvall, Torgny Lectures on the coupling method, Wiley Series in Probability and Mathematical Statistics, John Wiley \& Sons, Inc., New York, 1992 \href{https://ahl.centre-mersenne.org/item/AHL_2018__1__149_0/#Lin92}{Publisher reference}.
\bibitem[Mie01]{Mie01} Miermont, Grégory Ordered additive coalescent and fragmentations associated to Levy processes with no positive jumps, Electron. J. Probab., Volume 6 (2001) (art. no. 14, 33 p.) \href{https://ahl.centre-mersenne.org/item/AHL_2018__1__149_0/#Mie01}{Publisher reference}.
\bibitem[Mie08]{Mie08b} Miermont, Grégory Invariance principles for spatial multitype Galton-Watson trees, Ann. Inst. Henri Poincaré Probab. Stat., Volume 44 (2008) no. 6, pp. 1128-1161 \href{https://ahl.centre-mersenne.org/item/AHL_2018__1__149_0/#Mie08b}{Publisher reference}.
\bibitem[Mil77]{Mil77} Millar, P. W. Zero-one laws and the minimum of a Markov process, Trans. Amer. Math. Soc., Volume 226 (1977), pp. 365-391 \href{https://ahl.centre-mersenne.org/item/AHL_2018__1__149_0/#Mil77}{Publisher reference}.
\bibitem[MM03]{MM03} Marckert, Jean-François; Mokkadem, Abdelkader The depth first processes of Galton-Watson trees converge to the same Brownian excursion, Ann. Probab., Volume 31 (2003) no. 3, pp. 1655-1678 \href{https://ahl.centre-mersenne.org/item/AHL_2018__1__149_0/#MM03}{Publisher reference}.
\bibitem[MM07]{MM07} Marckert, Jean-François; Miermont, Grégory Invariance principles for random bipartite planar maps, Ann. Probab., Volume 35 (2007) no. 5, pp. 1642-1705 \href{https://ahl.centre-mersenne.org/item/AHL_2018__1__149_0/#MM07}{Publisher reference}.
\bibitem[Mos89]{Mos89} Moszkowski, Paul A solution to a problem of Dénes: a bijection between trees and factorizations of cyclic permutations, European J. Combin., Volume 10 (1989) no. 1, pp. 13-16 \href{https://ahl.centre-mersenne.org/item/AHL_2018__1__149_0/#Mos89}{Publisher reference}.
\bibitem[Nev86]{Nev86} Neveu, Jacques Arbres et processus de Galton-Watson, Ann. Inst. Henri Poincaré Probab. Stat., Volume 22 (1986) no. 2, pp. 199-207 \href{https://ahl.centre-mersenne.org/item/AHL_2018__1__149_0/#Nev86}{Publisher reference}.
\bibitem[NS06]{NS06} Nica, Alexandru; Speicher, Roland Lectures on the combinatorics of free probability, London Mathematical Society Lecture Note Series, 335, Cambridge University Press, Cambridge, 2006 \href{https://ahl.centre-mersenne.org/item/AHL_2018__1__149_0/#NS06}{Publisher reference}.
\bibitem[Oko00]{Oko00} Okounkov, Andrei Toda equations for Hurwitz numbers, Math. Res. Lett., Volume 7 (2000) no. 4, pp. 447-453 \href{https://ahl.centre-mersenne.org/item/AHL_2018__1__149_0/#Oko00}{Publisher reference}.
\bibitem[Pak99]{Pak1999} Pak, Igor Reduced decompositions of permutations in terms of star transpositions, generalized Catalan numbers and k-ary trees, Discrete Math., Volume 204 (1999), pp. 329-335 \href{https://ahl.centre-mersenne.org/item/AHL_2018__1__149_0/#Pak1999}{Publisher reference}.
\bibitem[Pit06]{Pit06} Pitman, Jim Combinatorial stochastic processes, Lecture Notes in Mathematics, 1875, Springer-Verlag, Berlin, 2006 (Lectures from the 32nd Summer School on Probability Theory held in Saint-Flour, July 7–24. With a foreword by Jean Picard) \href{https://ahl.centre-mersenne.org/item/AHL_2018__1__149_0/#Pit06}{Publisher reference}.
\bibitem[Sch05]{Scr05} Schramm, Oded Compositions of random transpositions, Israel J. Math., Volume 147 (2005), pp. 221-243 \href{https://ahl.centre-mersenne.org/item/AHL_2018__1__149_0/#Scr05}{Publisher reference}.
\bibitem[SSV97]{SSV97} Shapiro, Boris; Shapiro, Michael; Vainshtein, Alek Ramified coverings of s 2  with one degenerate branching point and enumeration of edge-ordered graphs, Adv. in Math. Sci., Volume 34 (1997), pp. 219-228 \href{https://ahl.centre-mersenne.org/item/AHL_2018__1__149_0/#SSV97}{Publisher reference}.
\bibitem[Sta84]{Sta84} Stanley, Richard P. On the number of reduced decompositions of elements of Coxeter groups, European J. Combin., Volume 5 (1984) no. 4, pp. 359-372 \href{https://ahl.centre-mersenne.org/item/AHL_2018__1__149_0/#Sta84}{Publisher reference}.
\bibitem[The]{Thev18} Thevenin, Paul (work in progress) \href{https://ahl.centre-mersenne.org/item/AHL_2018__1__149_0/#Thev18}{Publisher reference}.
\bibitem[Tót93]{Tot93} Tóth, Bálint Improved lower bound on the thermodynamic pressure of the spin 1/2 Heisenberg ferromagnet, Lett. Math. Phys., Volume 28 (1993) no. 1, pp. 75-84 \href{https://ahl.centre-mersenne.org/item/AHL_2018__1__149_0/#Tot93}{Publisher reference}.
\end{thebibliography}
\end{document}
